% Part: sets-functions-relations
% Chapter: relations
% Section: relations-as-sets

\documentclass[../../../include/open-logic-section]{subfiles}

\begin{document}

\olfileid{sfr}{rel}{set}
\olsection{سیٹاں دے روپ وچ تعلق}

\begin{explain}
\olref[sfr][set][imp]{sec} وچ اسی کجھ اہم سیٹاں دا ذکر کیتا سی:
$\Nat$, $\Int$, $\Rat$, $\Real$۔ تہانوں ضرور یاد ہووے گا کہ
ایہناں وچوں کجھ سیٹاں دے \pnboblique{!!{element}s} وچکار کجھ دلچسپ تعلق ہوندے نیں۔
مثال دے طور تے، ایہناں وچوں ہر سیٹ اُتے اک بالکل معیاری \emph{ترتیبی
تعلق} ہوندا اے۔ اک \emph{اکو ہون دا} تعلق وی اے، جیہڑا ہر شے دا
اپنے آپ نال ہوندا اے تے ہور کسے شے نال نہیں ہوندا۔ اگے سانوں
ہور کئی دلچسپ تعلق ملن گے، تے ممکن تعلق تاں اوہناں توں وی ودھ نیں۔
پر اوہناں دا جائزہ لین توں پہلاں اسی ایس گل توں شروع کراں گے کہ
تعلقاں نوں اک خاص قسم دیاں سیٹاں دے روپ وچ ویکھیا جا سکدا اے۔
  
ایس لئی \olref[sfr][set][pai]{sec} دیاں دو گلاں یاد کرو۔ پہلی،
\emph{ترتیب وار جوڑے} دا تصور یاد کرو: جے $a$ تے $b$ دتے ہون، تاں اسی
$\tuple{a, b}$ بنا سکدے آں۔ اہم گل ایہ اے کہ ایتھے عنصراں دی ترتیب
\emph{واقعی} فرق پاؤندی اے۔ سو جے $a \neq b$، تاں $\tuple{a, b} \neq \tuple{b, a}$۔
(ایہنوں بے ترتیب جوڑیاں، یعنی $2$ عنصراں والیاں سیٹاں، نال ملا کے ویکھو، جتھے
$\{a, b\}=\{b, a\}$۔) دوجی، \emph{کارتیسی
حاصلِ ضرب} دا تصور یاد کرو: جے $A$ تے $B$ سیٹ ہون، تاں اسی $A \times B$ بنا
سکدے آں؛ ایہ اوہناں سارے جوڑیاں $\tuple{x, y}$ دا سیٹ اے جنہاں لئی $x \in A$ تے
$y \in B$ ہووے۔ خاص طور تے، $A^{2}= A \times A$، $A$ توں بنن والے سارے ترتیب وار
جوڑیاں دا سیٹ اے۔
  
ہن اسی اک سیٹ اُتے اک خاص تعلق ویکھاں گے: $<$ دا تعلق، جیہڑا قدرتی عدداں
دے سیٹ $\Nat$ اُتے اے۔ عدداں دے اوہ سارے جوڑے $\tuple{n, m}$ لو جتھے $n<m$، یعنی
\[
  R=\Setabs{\tuple{n, m}}{n, m \in \Nat \text{ تے } n<m}.
\]
$n$ دا $m$ توں چھوٹا ہونا تے جوڑے $\tuple{n, m}$ دا $R$ دا رکن ہونا،
ایہناں وچکار ڈونگھا جوڑ اے، یعنی:
\[
      n<m\text{ اُدوں تے صرف اُدوں جدوں }\tuple{n, m} \in R.
\]
درحقیقت، معلومات دا کوئی حصہ گوائے بغیر اسی سیٹ $R$ نوں ای
$<$ دا تعلق \emph{منّ} سکدے آں، جیہڑا $\Nat$ اُتے اے۔

ایسے طرح عدداں وچکار کسے وی تعلق لئی اسی $\Nat^{2}$ دا اک ذیلی سیٹ بنا
سکدے آں۔ پُٹھے پاسے، جے عدداں دے جوڑیاں دا کوئی وی سیٹ
$S \subseteq \Nat^{2}$ دتا ہووے، تاں اوہدے مطابق عدداں وچکار اک تعلق ہوندا اے:
اوہ تعلق جیہڑا $n$ دا $m$ نال اُدوں تے صرف اُدوں ہووے جدوں
$\tuple{n, m} \in S$۔ ایس بنا اُتے اگلی تعریف جائز اے:
\end{explain}
  
\begin{defn}[دو رکنی تعلق]
اک سیٹ $A$ اُتے \emph{دو رکنی تعلق}، $A^{2}$ دا اک ذیلی سیٹ ہوندا اے۔ جے $R
\subseteq A^{2}$، $A$ اُتے اک دو رکنی تعلق ہووے تے $x, y \in A$ ہون، تاں
اسی کدی کدی $Rxy$ (یا $xRy$)، $\tuple{x, y} \in R$ دی تھاں لکھدے آں۔
\end{defn}
  
\begin{ex}
  \ollabel{relations}
قدرتی عدداں دے جوڑیاں دا سیٹ $\Nat^{2}$، دو جہتی میٹرکس وچ ایس طرح
لکھیا جا سکدا اے:
\[
  \begin{array}{ccccc}
  \mathbf{\tuple{ 0,0 }} & \tuple{ 0,1 } &
    \tuple{ 0,2 } & \tuple{ 0,3 } & \ldots\\
  \tuple{ 1,0 } & \mathbf{\tuple{ 1,1 }} &
    \tuple{ 1,2 } & \tuple{ 1,3 } & \ldots\\
  \tuple{ 2,0 } & \tuple{ 2,1 } &
    \mathbf{\tuple{ 2,2 }} & \tuple{ 2,3 } & \ldots\\
  \tuple{ 3,0 } & \tuple{ 3,1 } & \tuple{ 3,2 } &
    \mathbf{\tuple{ 3,3 }} & \ldots\\
  \vdots & \vdots & \vdots & \vdots & \mathbf{\ddots}
  \end{array}
\]
ایتھے اسی ترچھے قطر نوں موٹے اکھراں وچ وکھایا اے، کیوں جو $\Nat^2$ دا
اوہ ذیلی سیٹ جیہدے جوڑے ایس قطر اُتے پئے نیں، یعنی
\[
  \{\tuple{0,0 }, \tuple{ 1,1 }, \tuple{ 2,2 }, \dots\},
  \]
$\Nat$ اُتے \emph{اکو ہون دا تعلق} اے۔ (اکو ہون دا تعلق بہت ورتیا جاندا اے،
سو کسے وی سیٹ لئی $\Id{A}=\Setabs{\tuple{ x,x }}{x \in
A}$ دی تعریف کر لئیے، جتھے سیٹ $A$ اے۔) قطر توں اُتے پئے سارے جوڑیاں دا
ذیلی سیٹ، یعنی
\[
  L = \{\tuple{ 0,1 },\tuple{ 0,2 },\ldots,\tuple{ 1,2 },
  \tuple{ 1,3 }, \dots, \tuple{ 2,3 }, \tuple{ 2,4 },\ldots\},
\]
\emph{توں چھوٹا ہون} دا تعلق اے، یعنی $Lnm$ اُدوں تے صرف اُدوں جدوں $n<m$۔
قطر توں تھلے پئے جوڑیاں دا ذیلی سیٹ، یعنی
\[
  G=\{ \tuple{ 1,0 },\tuple{ 2,0 },\tuple{
    2,1 }, \tuple{ 3,0 },\tuple{ 3,1 },\tuple{ 3,2 }, \dots\},
\]
\emph{توں وڈا ہون} دا تعلق اے، یعنی $Gnm$ اُدوں تے صرف اُدوں جدوں $n>m$۔
$L$ دا $I$ نال اتحاد، جیہنوں اسی $K=L\cup I$ آکھ سکدے آں، \emph{توں چھوٹا
یا برابر ہون} دا تعلق اے: $Knm$ اُدوں تے صرف اُدوں جدوں $n \le m$۔ ایسے طرح
$H=G \cup I$، \emph{توں وڈا یا برابر ہون دا تعلق} اے۔ ایہ تعلق
$L$, $G$, $K$ تے $H$ اک خاص قسم دے تعلق نیں جنہاں نوں
\emph{ترتیباں} آکھدے نیں۔ $L$ تے $G$ دی ایہ خاصیت اے کہ کسے وی عدد دا اپنے
آپ نال نہ $L$ دا تعلق اے تے نہ $G$ دا (یعنی ہر $n$ لئی، نہ $Lnn$ تے نہ
$Gnn$)۔ ایس خاصیت والے تعلقاں نوں \emph{اپنے آپ نال تعلق نہ ہون والے} تعلق
آکھدے نیں؛ تے جے اوہ ترتیباں وی ہون، تاں اوہناں نوں \emph{سخت ترتیباں}
آکھدے نیں۔
\end{ex}

\begin{explain}
بھاویں ترتیباں تے اکو ہون دا تعلق اہم تے قدرتی تعلق نیں، پر ایس گل اُتے
زور دینا چاہیدا اے کہ ساڈی تعریف مطابق $A^{2}$ دا \emph{کوئی وی} ذیلی سیٹ،
$A$ اُتے اک تعلق اے، بھاویں اوہ کِنّا وی غیر فطری یا گھڑیا ہویا لگے۔ خاص
طور تے، $\emptyset$ کسے وی سیٹ اُتے اک تعلق اے (\emph{خالی تعلق}، جیہڑا
عنصراں دے کسے وی جوڑے وچکار نہیں ہوندا)، تے خود $A^{2}$ وی $A$ اُتے اک
تعلق اے (جیہڑا ہر جوڑے وچکار ہوندا اے)، جیہنوں \emph{ہمہ گیر تعلق} آکھدے
نیں۔ پر $E=\Setabs{\tuple{n, m}}{n>5 \text{ یا } m \times n \ge 34}$ ورگی
شے وی اک تعلق گنی جاندی اے۔
\end{explain}
  
\begin{prob}
  سیٹ اُتے تعلق $\subseteq$ دے !!{element}s دی فہرست بناؤ، جتھے سیٹ
  $\Pow{\{a, b, c\}}$ اے۔
\end{prob}

\end{document}
